\documentclass[a4paper,12pt]{article}

% -------------------------------------------------
% Кодировка и язык
% -------------------------------------------------
\usepackage[T2A]{fontenc}
\usepackage[utf8]{inputenc}
\usepackage[russian]{babel}

% -------------------------------------------------
% Математика
% -------------------------------------------------
\usepackage{amsmath,amssymb,amsthm,mathtools}
\usepackage{icomma}

% -------------------------------------------------
% Шрифт и типографика
% -------------------------------------------------
\usepackage{helvet}
\renewcommand{\familydefault}{\sfdefault}
\usepackage{microtype}
\usepackage{setspace}
\onehalfspacing

% -------------------------------------------------
% Геометрия страницы
% -------------------------------------------------
\usepackage[
  left=20mm,
  right=20mm,
  top=20mm,
  bottom=22mm
]{geometry}

% -------------------------------------------------
% Графика
% -------------------------------------------------
\usepackage{tikz}
\usetikzlibrary{calc}
\usetikzlibrary{arrows.meta,positioning,shapes.geometric}

% -------------------------------------------------
% Таблицы и списки
% -------------------------------------------------
\usepackage{tabularx,booktabs,longtable}
\usepackage{enumitem}

% -------------------------------------------------
% Служебные пакеты
% -------------------------------------------------
\usepackage{xcolor}
\usepackage{parskip}
\usepackage{fancyhdr}
\usepackage{hyperref}

\hypersetup{
  hidelinks,
  pdfauthor={Лёвин Артём Александрович},
  pdftitle={Степени. Формулы и стратегия преобразований}
}

% -------------------------------------------------
% Палитра
% -------------------------------------------------
\definecolor{accent}{HTML}{008080}
\definecolor{darkaccent}{HTML}{004D4D}
\definecolor{textgray}{HTML}{333333}
\definecolor{softgray}{HTML}{6A7373}

\color{textgray}

% -------------------------------------------------
% Колонтитулы
% -------------------------------------------------
\pagestyle{fancy}
\fancyhf{}
\fancyhead[L]{\small Дарья Савенкова}
\fancyhead[R]{\small ЕГЭ по математике}
\fancyfoot[C]{\small\thepage}
\renewcommand{\headrulewidth}{0.4pt}
\renewcommand{\headrule}{%
  \hbox to\headwidth{\color{accent}\leaders\hrule height \headrulewidth\hfill}
}

% -------------------------------------------------
% Настройки списков
% -------------------------------------------------
\setlist[itemize]{
  leftmargin=1.5em,
  itemsep=0.35em,
  topsep=0.4em
}

\setlist[enumerate]{
  leftmargin=1.8em,
  itemsep=0.55em,
  topsep=0.5em
}

\setlength{\parindent}{0pt}
\emergencystretch=2em

% -------------------------------------------------
% Стили блок-схем
% -------------------------------------------------
\tikzset{
  algstart/.style={
    ellipse,
    draw=darkaccent,
    line width=0.9pt,
    align=center,
    minimum width=5.1cm,
    minimum height=1.15cm,
    text width=4.6cm
  },
  algprocess/.style={
    rectangle,
    rounded corners=3pt,
    draw=accent,
    line width=0.9pt,
    align=center,
    minimum height=1.2cm,
    text width=5.1cm,
    inner sep=6pt
  },
  algside/.style={
    rectangle,
    rounded corners=3pt,
    draw=accent,
    line width=0.9pt,
    align=center,
    minimum height=1.15cm,
    text width=4.0cm,
    inner sep=5pt
  },
  algdecision/.style={
    diamond,
    aspect=2.25,
    draw=darkaccent,
    line width=0.9pt,
    align=center,
    text width=3.7cm,
    inner sep=1.5pt
  },
  algarrow/.style={
    -{Stealth[length=3mm,width=2mm]},
    line width=0.9pt,
    draw=darkaccent
  }
}

% -------------------------------------------------
% Вспомогательные команды
% -------------------------------------------------
\newcommand{\lessonrule}{%
  \par\vspace{0.3em}
  {\color{accent}\rule{\linewidth}{0.6pt}}
  \vspace{0.5em}
}

\newcommand{\important}[1]{\textbf{\color{darkaccent}#1}}

\begin{document}

% =================================================
% ТИТУЛЬНЫЙ ЛИСТ
% =================================================
\begin{titlepage}
\thispagestyle{empty}

\begin{center}
\vspace*{23mm}

{\Large\color{darkaccent}\textbf{ЧЕК-ЛИСТ}}

\vspace{8mm}

{\huge\bfseries
Степени}

\vspace{3mm}

{\Large
формулы и стратегия преобразований}

\vspace{15mm}

{\large
Подготовка к ЕГЭ по профильной математике}

\vfill

{\large
\textbf{Лёвин Артём Александрович}\\[2mm]
эксклюзивно для Дарьи Савенковой}

\vspace{11mm}

{\large \today}

\vspace{22mm}
\end{center}

% Сдержанная кривая Безье внизу титульного листа
\begin{tikzpicture}[remember picture,overlay]
  \draw[
    accent,
    line width=1.15pt
  ]
  ([xshift=18mm,yshift=17mm]current page.south west)
  .. controls
  ([xshift=55mm,yshift=27mm]current page.south west)
  and
  ([xshift=-60mm,yshift=7mm]current page.south east)
  ..
  ([xshift=-18mm,yshift=17mm]current page.south east);
\end{tikzpicture}

\end{titlepage}

% =================================================
% ОСНОВНОЙ ЧЕК-ЛИСТ
% =================================================
\section*{\color{darkaccent}Основной чек-лист}
\lessonrule

\subsection*{\color{accent}1. Что нужно видеть в любой степени}

В выражении
\[
a^n
\]
число или выражение \(a\) называется \important{основанием степени},
а \(n\) --- \important{показателем степени}.

Если показатель явно не записан, подразумевается единица:
\[
a=a^1.
\]

Особые случаи:
\[
a^1=a,
\]
\[
a^0=1,\qquad a\ne0.
\]

Следовательно,
\[
0^0
\]
в школьной алгебре не рассматривается как определённое значение.

\subsection*{\color{accent}2. Главная идея занятия}

При преобразовании выражений со степенями удобно разделять два уровня работы:

\begin{enumerate}
  \item \important{Принцип} подсказывает, с чего начать.
  \item \important{Формула} позволяет технически выполнить преобразование.
\end{enumerate}

Главный принцип:

\[
\boxed{
\text{составные числа}
\;\longrightarrow\;
\text{корни}
\;\longrightarrow\;
\text{дроби}
}
\]

Их по возможности приводим к удобному \important{степенному виду}.

Важно: этот принцип прежде всего относится к \important{основаниям степеней}.
Показатель степени сам по себе может быть дробным, десятичным,
содержать корень или другое выражение.

Например, в
\[
5^{0,36}\cdot25^{0,32}
\]
преобразовать требуется число \(25\), поскольку это основание:
\[
25=5^2.
\]

Показатели \(0,36\) и \(0,32\) уже являются показателями степеней,
поэтому специально преобразовывать их перед началом решения не требуется.

% =================================================
% ФОРМУЛЫ
% =================================================
\section*{\color{darkaccent}Формулы степеней}
\lessonrule

Все формулы ниже используются только там, где соответствующие выражения
определены.

При работе с произвольными дробными или действительными показателями
самый надёжный универсальный режим --- \important{положительные основания}.

\renewcommand{\arraystretch}{1.45}

\begin{table}[ht]
\centering
\caption{Основные свойства степеней}
\begin{tabularx}{\linewidth}{@{}>{\bfseries}c X X@{}}
\toprule
№ & Формула & Условие и смысл \\
\midrule

1 &
\[
a^m\cdot a^n=a^{m+n}
\]
&
Одинаковые основания и умножение:
показатели складываем.
\\

2 &
\[
\frac{a^m}{a^n}=a^{m-n}
\]
&
Одинаковые основания и деление,
\(a\ne0\):
показатели вычитаем.
\\

3 &
\[
(ab)^n=a^n b^n
\]
&
Степень произведения распределяется
на каждый множитель.
\\

4 &
\[
\left(\frac{a}{b}\right)^n
=
\frac{a^n}{b^n}
\]
&
Степень частного распределяется
на числитель и знаменатель,
\(b\ne0\).
\\

5 &
\[
(a^m)^n=a^{mn}
\]
&
При возведении степени в степень
показатели перемножаем.
\\

6 &
\[
\frac{1}{a^n}=a^{-n}
\]
&
Отрицательный показатель позволяет
перейти от дроби к степени,
\(a\ne0\).
\\

7 &
\[
\left(\frac{a}{b}\right)^n
=
\left(\frac{b}{a}\right)^{-n}
\]
&
При обращении дроби знак показателя
меняется,
\(a\ne0,\ b\ne0\).
\\

8 &
\[
\sqrt[k]{a^n}=a^{\frac nk}
\]
&
Надёжно применяется при \(a>0\).
Внутренний показатель делим
на показатель корня.
\\

\bottomrule
\end{tabularx}
\end{table}

\subsection*{\color{accent}Формулы работают в обе стороны}

Полезно уметь узнавать каждую формулу как слева направо,
так и справа налево.

Например:
\[
a^m\cdot a^n=a^{m+n}
\]
можно использовать в обратную сторону:
\[
a^{m+n}=a^m\cdot a^n.
\]

А формула
\[
a^n b^n=(ab)^n
\]
позволяет \important{вынести общий показатель степени}.

Именно поэтому при решении задачи полезно постоянно задавать себе вопросы:

\begin{itemize}
  \item можно ли раскрыть скобки;
  \item можно ли собрать одинаковые основания;
  \item можно ли вынести общий показатель;
  \item появилась ли возможность применить следующую формулу.
\end{itemize}

% =================================================
% ГРАНИЦЫ ПРИМЕНИМОСТИ
% =================================================
\section*{\color{darkaccent}Границы применимости и типичные ошибки}
\lessonrule

\subsection*{\color{accent}1. Одинаковое основание обязательно}

Можно:
\[
5^2\cdot5^4=5^6.
\]

Нельзя складывать показатели при разных основаниях:
\[
5^2\cdot3^4
\ne
15^6.
\]

\subsection*{\color{accent}2. Умножение и сложение --- разные операции}

Верно:
\[
5^2\cdot5^4=5^6.
\]

Но
\[
5^2+5^4
\ne
5^6.
\]

Наличие знаков \(+\) и \(-\) внутри выражения --- сигнал
особенно внимательно проверить возможность применения формулы.

В частности,
\[
(a+b)^n\ne a^n+b^n
\]
в общем случае.

Например:
\[
(2+3)^2=25,
\]
тогда как
\[
2^2+3^2=13.
\]

\subsection*{\color{accent}3. Не путайте две похожие ситуации}

Если
\[
5^3\cdot5^4,
\]
то показатели \important{складываются}:
\[
5^{3+4}=5^7.
\]

Если
\[
(5^3)^4,
\]
то показатели \important{перемножаются}:
\[
5^{3\cdot4}=5^{12}.
\]

\subsection*{\color{accent}4. Отрицательная степень}

При \(a\ne0\):
\[
a^{-n}=\frac1{a^n}.
\]

Например:
\[
2^{-1}=\frac12,
\qquad
5^{-3}=\frac1{125}.
\]

Отрицательная степень не означает отрицательное число.

Например:
\[
2^{-3}=\frac18>0.
\]

\subsection*{\color{accent}5. Корни и дробные показатели}

Для положительного \(a\):
\[
\sqrt[k]{a^n}=a^{\frac nk}.
\]

Например:
\[
\sqrt[7]{3^5}=3^{\frac57}.
\]

Если перед числом показатель отсутствует:
\[
5=5^1,
\]
поэтому
\[
\sqrt{5}=5^{\frac12}.
\]

Для корня чётной степени в вещественных числах требуется
неотрицательный подкоренный объект:
\[
\sqrt[2k]{A}
\quad\Longrightarrow\quad
A\ge0.
\]

Если такой корень находится в знаменателе, дополнительно требуется
\[
A>0.
\]

При преобразованиях дробных степеней с отрицательными основаниями
необходимо отдельно контролировать область определения.
В задачах ЕГЭ при возможности безопаснее сначала явно установить ОДЗ.

% =================================================
% СКОБКИ
% =================================================
\section*{\color{darkaccent}Работа со скобками}
\lessonrule

Скобки --- один из главных технических инструментов алгебры.

\subsection*{\color{accent}Если скобки появились, ищите возможность их раскрыть}

Например:
\[
(5^2)^4=5^8,
\]
\[
(2\cdot3)^n=2^n3^n.
\]

\subsection*{\color{accent}Если повторяются элементы, ищите возможность собрать их}

Например:
\[
a^n b^n=(ab)^n.
\]

\subsection*{\color{accent}Мнемоника «ВУПС»}

Полезное правило:

\[
\boxed{\text{ВУПС}}
\]

\begin{itemize}
  \item \textbf{В} --- вычитаемое;
  \item \textbf{У} --- умножаемое;
  \item \textbf{П} --- подставляемое;
  \item \textbf{С} --- берём в скобки.
\end{itemize}

Особенно важны скобки, если внутри выражения уже имеются знаки
\(+\) или \(-\).

Например:
\[
A-\bigl(2\sqrt7-1\bigr).
\]

После раскрытия:
\[
A-2\sqrt7+1.
\]

Если забыть скобки, знак перед вторым слагаемым легко потерять.

% =================================================
% СОСТАВНЫЕ ЧИСЛА
% =================================================
\section*{\color{darkaccent}Составные числа и факторизация}
\lessonrule

\subsection*{\color{accent}Простые и составные числа}

\important{Простое число} имеет ровно два натуральных делителя:
\(1\) и само себя.

Например:
\[
2,\ 3,\ 5,\ 7,\ 11,\ 13.
\]

\important{Составное число} имеет более двух натуральных делителей.

Например, число \(10\) имеет делители
\[
1,\ 2,\ 5,\ 10,
\]
поэтому оно составное.

В задачах со степенями составное число часто требуется
представить как степень или произведение простых множителей.

Например:
\[
8=2^3,
\]
\[
25=5^2,
\]
\[
27=3^3,
\]
\[
81=3^4.
\]

\subsection*{\color{accent}Факторизация}

Если подходящая степень не видна сразу, число раскладываем
на простые множители.

Для \(81\):
\[
81=3\cdot27
=3\cdot3\cdot9
=3\cdot3\cdot3\cdot3
=3^4.
\]

Следовательно,
\[
81=3^4.
\]

Этот приём особенно полезен при работе с:

\begin{itemize}
  \item корнями;
  \item степенями;
  \item логарифмами;
  \item квадратными уравнениями;
  \item выражениями повышенной сложности.
\end{itemize}

% =================================================
% ПРИНЦИП ПРИВЕДЕНИЯ К СТЕПЕННОМУ ВИДУ
% =================================================
\section*{\color{darkaccent}Принцип: приводим к степенному виду}
\lessonrule

Получив выражение со степенями, сначала просматриваем
\important{основания}.

\subsection*{\color{accent}Шаг 1. Составные числа}

Например:
\[
25=5^2.
\]

\subsection*{\color{accent}Шаг 2. Корни}

Например:
\[
\sqrt{8}
=
\sqrt{2^3}
=
2^{\frac32}.
\]

Обратите внимание на последовательность:

\[
8
\longrightarrow
2^3
\longrightarrow
2^{\frac32}.
\]

Сначала обработано составное число,
затем применена формула для корня.

\subsection*{\color{accent}Шаг 3. Дроби}

Например:
\[
\frac18
=
\frac1{2^3}
=
2^{-3}.
\]

И снова сначала:
\[
8=2^3,
\]
после чего используется отрицательная степень.



% =================================================
% ПРИМЕР 1
% =================================================
\section*{\color{darkaccent}Разобранные примеры}
\lessonrule

\subsection*{\color{accent}Пример 1. Составное основание}

Вычислим:
\[
5^{0,36}\cdot25^{0,32}.
\]

Число \(25\) является составным основанием:
\[
25=5^2.
\]

Подставляем:
\[
5^{0,36}\cdot(5^2)^{0,32}.
\]

Раскрываем степень в степени:
\[
(5^2)^{0,32}
=
5^{2\cdot0,32}
=
5^{0,64}.
\]

Получаем:
\[
5^{0,36}\cdot5^{0,64}
=
5^{0,36+0,64}
=
5^1
=
\boxed{5}.
\]

\important{Что было главным:}
первое действие появилось из принципа
«составное основание приводим к степенному виду».

% =================================================
% ПРИМЕР 2
% =================================================
\subsection*{\color{accent}Пример 2. Составное число в знаменателе}

Вычислим:
\[
\frac{2^{\frac72}\cdot3^{\frac{11}{2}}}
     {6^{\frac92}}.
\]

Составное число:
\[
6=2\cdot3.
\]

Подставляем:
\[
\frac{2^{\frac72}\cdot3^{\frac{11}{2}}}
     {(2\cdot3)^{\frac92}}.
\]

Раскрываем степень произведения:
\[
(2\cdot3)^{\frac92}
=
2^{\frac92}3^{\frac92}.
\]

Тогда
\[
\frac{2^{\frac72}3^{\frac{11}{2}}}
     {2^{\frac92}3^{\frac92}}
=
2^{\frac72-\frac92}
\cdot
3^{\frac{11}{2}-\frac92}.
\]

Получаем:
\[
2^{-1}\cdot3
=
\frac12\cdot3
=
\boxed{\frac32}.
\]

Здесь отрицательная степень стала промежуточным,
полностью нормальным результатом.

% =================================================
% ПРИМЕР 3
% =================================================
\subsection*{\color{accent}Пример 3. Корень переводим в степень}

Вычислим:
\[
\left(
\frac{
2^{\frac13}\cdot2^{\frac14}
}{
\sqrt[12]{2}
}
\right)^2.
\]

Корень представляем степенью:
\[
\sqrt[12]{2}=2^{\frac1{12}}.
\]

Тогда
\[
\left(
2^{\frac13+\frac14-\frac1{12}}
\right)^2.
\]

Приводим показатели к общему знаменателю:
\[
\frac13+\frac14-\frac1{12}
=
\frac4{12}+\frac3{12}-\frac1{12}
=
\frac6{12}
=
\frac12.
\]

Следовательно,
\[
\left(2^{\frac12}\right)^2
=
2^1
=
\boxed{2}.
\]

% =================================================
% ПРИМЕР 4
% =================================================
\subsection*{\color{accent}Пример 4. Выражения в показателях и правило скобок}

Вычислим:
\[
\frac{
5^{3\sqrt7-1}\cdot5^{1-\sqrt7}
}{
5^{2\sqrt7-1}
}.
\]

Основания одинаковые, поэтому работаем с показателями:
\[
5^{
(3\sqrt7-1)
+
(1-\sqrt7)
-
(2\sqrt7-1)
}.
\]

Вычитаемый показатель обязательно взят в скобки.

Раскрываем:
\[
3\sqrt7-1+1-\sqrt7-2\sqrt7+1.
\]

Собираем подобные:
\[
3\sqrt7-\sqrt7-2\sqrt7=0,
\]
\[
-1+1+1=1.
\]

Следовательно,
\[
5^1=\boxed{5}.
\]

% =================================================
% ДВОЕТОЧИЕ
% =================================================
\section*{\color{darkaccent}Технический приём: деление записываем дробью}
\lessonrule

Если выражение записано через двоеточие,
часто полезно сразу перейти к дробной черте:
\[
A:B=\frac AB,
\qquad B\ne0.
\]

Например:
\[
5^x:5^y
=
\frac{5^x}{5^y}
=
5^{x-y}.
\]

Такая запись делает структуру выражения нагляднее
и снижает риск ошибки со знаками.

% =================================================
% САМОПРОВЕРКА
% =================================================
\section*{\color{darkaccent}Финальная самопроверка}
\lessonrule

Перед завершением решения проверьте:

\begin{enumerate}
  \item Определены ли все исходные выражения?
  \item Нет ли нуля в знаменателе?
  \item Выполнены ли ограничения для корней чётной степени?
  \item Проверены ли основания с отрицательными и дробными показателями?
  \item Все ли составные основания приведены к удобному виду?
  \item Все ли корни, которые мешали работе, представлены степенями?
  \item Все ли дроби при необходимости преобразованы через отрицательные степени?
  \item Не применено ли сложение показателей там, где между степенями стоит \(+\)?
  \item Не перепутаны ли правила
  \[
  a^m a^n=a^{m+n}
  \quad\text{и}\quad
  (a^m)^n=a^{mn}?
  \]
  \item Взяты ли в скобки сложные вычитаемые, умножаемые и подставляемые выражения?
  \item После каждого преобразования появилась ли более простая структура
  или новая возможность применить формулу?
  \item Проверена ли итоговая арифметика?
\end{enumerate}

% =================================================
% ТРЕНИРОВОЧНЫЙ БЛОК
% =================================================
\clearpage

\section*{\color{darkaccent}Тренировочный блок: 10 задач}
\lessonrule

Задачи расположены по нарастающей сложности.
Решения выполняйте последовательно, стараясь каждый раз сначала определить
\important{принцип}, а затем подобрать нужную \important{формулу}.

\subsection*{\color{accent}Средний уровень}

\begin{enumerate}[label=\textbf{\arabic*.}]

\item Вычислите:
\[
\frac{2^7\cdot2^3}{2^5}.
\]

\item Вычислите:
\[
\frac{(3^2)^4}{3^5}.
\]

\item Представьте
\[
\frac1{5^3}
\]
в виде степени с основанием \(5\).

\end{enumerate}

\subsection*{\color{accent}Уровень выше среднего}

\begin{enumerate}[label=\textbf{\arabic*.},start=4]

\item Вычислите:
\[
16^{\frac34}.
\]

\item Вычислите:
\[
27^{\frac23}.
\]

\item Вычислите:
\[
5^{0,38}\cdot25^{0,31}.
\]

\item Вычислите:
\[
\frac{
2^{\frac72}\cdot3^{\frac{11}{2}}
}{
6^{\frac92}
}.
\]

\item Вычислите:
\[
\left(
\frac{
2^{\frac13}\cdot2^{\frac14}
}{
\sqrt[12]{2}
}
\right)^2.
\]

\item Вычислите:
\[
\left(\frac23\right)^{-4}
\cdot
\left(\frac32\right)^{-2}.
\]

\end{enumerate}

\subsection*{\color{accent}Сложная задача}

\begin{enumerate}[label=\textbf{\arabic*.},start=10]

\item Вычислите:
\[
\frac{
5^{3\sqrt7-1}
\cdot
25^{1-\sqrt7}
}{
125^{\frac{\sqrt7}{3}}
}.
\]

\end{enumerate}

% =================================================
% ОТВЕТЫ
% =================================================
\clearpage

\section*{\color{darkaccent}Ответы к тренировочному блоку}
\lessonrule

\begin{table}[ht]
\centering
\caption{Таблица ответов}
\renewcommand{\arraystretch}{1.4}
\begin{tabularx}{0.72\linewidth}{@{}>{\bfseries}c X@{}}
\toprule
№ & Ответ \\
\midrule
1 & \(32\) \\
2 & \(27\) \\
3 & \(5^{-3}\) \\
4 & \(8\) \\
5 & \(9\) \\
6 & \(5\) \\
7 & \(\dfrac32\) \\
8 & \(2\) \\
9 & \(\dfrac94\) \\
10 & \(5\) \\
\bottomrule
\end{tabularx}
\end{table}

\vfill

\begin{center}
{\color{softgray}\small
Ключевая стратегия:
сначала привести основания к удобному степенному виду,
затем последовательно применять свойства степеней,
контролируя ОДЗ и знаки.}
\end{center}

\end{document}